\section{Our Learning Reflection}

From this analysis, we learned that statistics is not only about putting
numbers into formulas, but also about understanding and explaining the
meaning behind the results.

The first step is understanding the data, such as the number of
respondents, the variables being studied, and how the data were collected.
Descriptive statistics then helps us understand the characteristics of the
data through values such as mean, median, and standard deviation.
Visualization also helps us identify patterns in the data more easily.

We also learned that bivariate analysis can be used to examine the
relationship between two variables. In this research, the Pearson
correlation produced a value of $r = \resR$, indicating a strong positive
relationship between variables X and Y. In addition, the $R^2$ value of
\resRSquaredPct\% helped us understand how much variation in Y can be
explained by X within the research model.

From the t-test and F-test, we learned that statistical significance is
also important when interpreting research results. The significance value
of $\resSig < 0.05$ indicates that the statistical tests produced
significant results within the tested model.

We also learned how to apply probability using the same research data.
With \resN{} respondents, we can define the sample space, create events A and
B, and calculate probability, intersection, and conditional probability.
For example, the probability of a respondent having a high Y score is
\resPB\%, while when we know that the respondent has a high X score, the
probability of having a high Y score becomes \resPBGivenA\%.

The most important thing we learned is interpretation. Statistical values
are not meaningful enough if they are only presented as numbers. They need
to be translated into information that can be understood.

We also learned that a statistically significant relationship does not
automatically prove a cause-and-effect relationship. Other factors may
also be related to changes in variable Y. In this research, the $R^2$
value of \resRSquaredPct\% indicates that approximately
\resUnexplainedPct\% of the variation in Y is not explained by X within
the model.

Through this process, we began to see statistics as a way to tell the
story of data logically, understand patterns, and draw conclusions based
on the available evidence.
